\section{Limitations}\label{sec:limits}

The limitations of the kit and of the case study are those of the short paper.

% written by scripts/reuse_short.py from the sources of the short paper; do not edit
\begin{itemize}
\tightlist
\item
  \textbf{Solve-level parity is not refuted, not established.} The evidence (Figure~\ref{fig:parity}A) is one 24-task list, one repeat per run; ten runs solved only six distinct tasks.
\item
  \textbf{Kaggle-side facts are partly inferred}: the hosted google-adk release, tool-schema equality, and the identity of the hosted template and weights with the public Kaggle model.
\item
  \textbf{Environment differences remain.} One RTX 5090 instead of four L4s; 2 CPUs and 10 GiB per job; a bare \texttt{pytest} command, used in 30 of 216 commands in our Kaggle traces, is absent locally.
\item
  \textbf{The case study} covers one model, one bundle, one google-adk release and one configuration per divergence. Arm T swaps a whole template; in arm H thoughts still reach compaction summaries.
\end{itemize}


The record of experiments adds the following.

\begin{itemize}[leftmargin=*,itemsep=2pt,topsep=3pt]
\item \textbf{Rows that are missing.} One 31B block, an arm with a ten-minute limit that solved 18 of 30 development
tasks, has no result rows that we could retrieve. It is in no table and in no released file. The rows of one
workstation parity run under runner version 2 (4 of 24 solved) were not kept either; that run appears only in
\cref{fig:parity}.
\item \textbf{One run per task for 31B.} On the development tasks every 31B configuration ran each task once.
Differences between configurations and differences between runs cannot be separated there (\cref{sec:big}).
\item \textbf{Arms and sessions are confounded in places.} Some 31B arms ran in blocks of their own, under a
different load. The E4B candidates always ran in one job queue with the baseline.
\item \textbf{Task sets were not always chosen by a recorded rule.} The tasks of the two 31B blocks on real hardware
were picked with earlier results in view. For the eight tasks of block \texttt{berthE4-k2} the method of selection
was not recorded. The screen set S is not fresh: 32 of its 48 tasks had been run in earlier blocks. The development
block of the third programme used 16 tasks of the small panel after that panel's confirmation.
\item \textbf{The hold-out H was used twice.} The third candidate was built after the result of the first
confirmation on H was known. Its confirmation was registered with this disclosure, and its verdict is ``not
confirmed''.
\item \textbf{Post hoc analyses.} The noise estimates, the selection effect of \cref{sec:leadgain} and all of
\cref{sec:analysis} were chosen after the results were known. They describe the record. They are not tests.
\item \textbf{One baseline.} The share of tasks with an uncertain outcome, and with it the run-to-run noise, was
measured for one E4B agent on Berth tasks. Another agent, another model or another task source can have a different
share.
\item \textbf{Hosted scores are single draws.} We have one hosted draw per file and no repeated draw of an E4B file.
\item \textbf{Machines differ between experiments.} The 31B runs on development tasks used emulated latency on a
single GPU, the first E4B blocks used hosted notebooks with four L4 GPUs, and most E4B runs used the workstation.
Times are comparable within a block and not across machines.
\item \textbf{Arms cannot be rebuilt from this paper.} Prompts, the E4B bundle and the adapter are withheld
(\cref{sec:setting}), and each arm is described by the kind of its change. The rows of the development tasks carry
opaque task codes, because the competition data may not be redistributed; they support counts and pairings and
cannot be joined with the task texts.
\end{itemize}
